\noindent\textbf{Elección del procedimiento.} Los grupos de similitud usan $\mathrm{Ar}=\rho Bd^3/\mu^2$. Al no existir ese Arquímedes en la tabla, se interpola $\ln U$ frente a $\ln B$ entre 20 y 40 N/m$^3$: $n=\ln(U_{40}/U_{20})/\ln2$ y $U_{30}=U_{20}(30/20)^n$. Richardson compara flotabilidad con inercia: $\mathrm{Ri}=Bd/(\rho U^2)$.

\begin{enumerate}

\item Identificar el punto adimensional que falta en la tabla.

\[\fitmath{\mathrm{Ar}=\frac{\rho\cdot B\cdot {d}^{3}}{{\mu}^{2}}}\]

\[\fitmath{\mathrm{Ar}=\frac{(1000)\cdot (30)\cdot {(0{,}004)}^{3}}{{(0{,}001)}^{2}}\approx 1920\,}\]

\item Calcular pendiente de la recta en ejes logarítmicos.

\[\fitmath{n=\frac{\ln\left(\frac{U_{40}}{U_{20}}\right)}{\ln\left(2\right)}}\]

\[\fitmath{n=\frac{\ln\left(\frac{(0{,}0527)}{(0{,}0428)}\right)}{\ln\left(2\right)}\approx 0{,}3001922\,}\]

\item Volver de la interpolación logarítmica a velocidad.

\[\fitmath{\mathrm{Uc}=U_{20}\cdot {\left(\frac{B}{B_{20}}\right)}^{n}}\]

\[\fitmath{\mathrm{Uc}=(0{,}0428)\cdot {\left(\frac{(30)}{(20)}\right)}^{(0{,}3001921652)}\approx 0{,}04833982\,\mathrm{m/s}}\]

\item Usar la velocidad interpolada sin redondear antes.

\[\fitmath{\mathrm{Ri}=\frac{B\cdot d}{\rho\cdot {\mathrm{Uc}}^{2}}}\]

\[\fitmath{\mathrm{Ri}=\frac{(30)\cdot (0{,}004)}{(1000)\cdot {(0{,}04833981515)}^{2}}\approx 0{,}05135365\,}\]

\item Para evitar interpolar, imponer $240d_{nuevo}^3=Bd^3$ con el mismo líquido.

\[\fitmath{\mathrm{dnew}=d\cdot {\left(\frac{B}{\mathrm{Bnew}}\right)}^{\frac{1}{3}}}\]

\[\fitmath{\mathrm{dnew}=(0{,}004)\cdot {\left(\frac{(30)}{(240)}\right)}^{\frac{1}{3}}\approx 0{,}002\,\mathrm{m}}\]

\end{enumerate}
